\documentclass{amsart}
% For dealing with references we use the comment environment
\usepackage{verbatim}
\newenvironment{reference}{\comment}{\endcomment}
%\newenvironment{reference}{}{}
\newenvironment{slogan}{\comment}{\endcomment}
\newenvironment{history}{\comment}{\endcomment}

% For commutative diagrams we use Xy-pic
\usepackage[all]{xy}

% We use 2cell for 2-commutative diagrams.
\xyoption{2cell}
\UseAllTwocells

% We use multicol for the list of chapters between chapters
\usepackage{multicol}

% This is generally recommended for better output
\usepackage{lmodern}
\usepackage[T1]{fontenc}

% For cross-file-references

% Package for hypertext links:
\usepackage{hyperref}

% For any local file, say "hello.tex" you want to link to please

% Theorem environments.
%
\theoremstyle{plain}
\newtheorem{theorem}[subsection]{Theorem}
\newtheorem{proposition}[subsection]{Proposition}
\newtheorem{lemma}[subsection]{Lemma}

\theoremstyle{definition}
\newtheorem{definition}[subsection]{Definition}
\newtheorem{example}[subsection]{Example}
\newtheorem{exercise}[subsection]{Exercise}
\newtheorem{situation}[subsection]{Situation}

\theoremstyle{remark}
\newtheorem{remark}[subsection]{Remark}
\newtheorem{remarks}[subsection]{Remarks}

\numberwithin{equation}{subsection}

% Macros
%
\def\lim{\mathop{\mathrm{lim}}\nolimits}
\def\colim{\mathop{\mathrm{colim}}\nolimits}
\def\Spec{\mathop{\mathrm{Spec}}}
\def\Hom{\mathop{\mathrm{Hom}}\nolimits}
\def\Ext{\mathop{\mathrm{Ext}}\nolimits}
\def\SheafHom{\mathop{\mathcal{H}\!\mathit{om}}\nolimits}
\def\SheafExt{\mathop{\mathcal{E}\!\mathit{xt}}\nolimits}
\def\Sch{\mathit{Sch}}
\def\Mor{\mathop{\mathrm{Mor}}\nolimits}
\def\Ob{\mathop{\mathrm{Ob}}\nolimits}
\def\Sh{\mathop{\mathit{Sh}}\nolimits}
\def\NL{\mathop{N\!L}\nolimits}
\def\CH{\mathop{\mathrm{CH}}\nolimits}
\def\proetale{{pro\text{-}\acute{e}tale}}
\def\etale{{\acute{e}tale}}
\def\QCoh{\mathit{QCoh}}
\def\Ker{\mathop{\mathrm{Ker}}}
\def\Im{\mathop{\mathrm{Im}}}
\def\Coker{\mathop{\mathrm{Coker}}}
\def\Coim{\mathop{\mathrm{Coim}}}

% Boxtimes
%
\DeclareMathSymbol{\boxtimes}{\mathbin}{AMSa}{"02}

%
% Macros for moduli stacks/spaces
%
\def\QCohstack{\mathcal{QC}\!\mathit{oh}}
\def\Cohstack{\mathcal{C}\!\mathit{oh}}
\def\Spacesstack{\mathcal{S}\!\mathit{paces}}
\def\Quotfunctor{\mathrm{Quot}}
\def\Hilbfunctor{\mathrm{Hilb}}
\def\Curvesstack{\mathcal{C}\!\mathit{urves}}
\def\Polarizedstack{\mathcal{P}\!\mathit{olarized}}
\def\Complexesstack{\mathcal{C}\!\mathit{omplexes}}
% \Pic is the operator that assigns to X its picard group, usage \Pic(X)
% \Picardstack_{X/B} denotes the Picard stack of X over B
% \Picardfunctor_{X/B} denotes the Picard functor of X over B
\def\Pic{\mathop{\mathrm{Pic}}\nolimits}
\def\Picardstack{\mathcal{P}\!\mathit{ic}}
\def\Picardfunctor{\mathrm{Pic}}
\def\Deformationcategory{\mathcal{D}\!\mathit{ef}}

\setlength{\emergencystretch}{3em}
\begin{document}
\title[Stacks Project, Tag 0BN4]{296 --- Classification of Galois categories by finite continuous automorphism-group sets}
\maketitle
\noindent Copyright (C) 2005--2025 Johan de Jong. Stacks Project authors. Source: \href{https://stacks.math.columbia.edu/tag/0BN4}{Tag 0BN4}.

\noindent Editorial difficulty note (not author text). Difficulty H2: The axioms must produce both full faithfulness and essential surjectivity rather than assume the category already comes from a profinite group. The proof realizes equivariant maps as graphs of connected components and constructs every coset object by a quotient of a Galois object whose stabilizer lies in the prescribed open subgroup..

\noindent Editorial provenance note (not author text). History: excerpt from revision \texttt{a0444\allowbreak 6e57e\allowbreak c1fbc\allowbreak 252a8\allowbreak 71afc\allowbreak ec775\allowbreak 2fb28\allowbreak 07b14}, pione.tex, lines 611--669. Reference presentation was adapted; original-author context and core support blocks were retained or added. The mathematical wording is unchanged. Licensed under GNU FDL 1.2 or later; no invariant sections or cover texts. See ../licenses/COPYING. Context blocks reproduce author definitions and supporting results; they are not additional counted problems.

\begin{definition}
\label{definition-galois-category}
\begin{reference}
Different from the definition in \href{https://stacks.math.columbia.edu/bibliography}{Bibliography: SGA1 (Expos\textbackslash{}'e V, Definition 5.1)}.
Compare with \href{https://stacks.math.columbia.edu/bibliography}{Bibliography: BS (Definition 7.2.1)}.
\end{reference}
Let $\mathcal{C}$ be a category and let $F : \mathcal{C} \to \textit{Sets}$
be a functor. The pair $(\mathcal{C}, F)$ is a {\it Galois category} if
\begin{enumerate}
\item $\mathcal{C}$ has finite limits and finite colimits,
\item
\label{item-connected-components}
every object of $\mathcal{C}$ is a finite (possibly empty)
coproduct of connected objects,
\item $F(X)$ is finite for all $X \in \Ob(\mathcal{C})$, and
\item $F$ reflects isomorphisms\footnote{Namely, given a morphism
$f$ of $\mathcal{C}$ if $F(f)$ is an isomorphism, then
$f$ is an isomorphism.} and is exact\footnote{This means that
$F$ commutes with finite limits and colimits, see
Categories, Section \href{https://stacks.math.columbia.edu/tag/0033}{section exact functor (Tag 0033)}.}.
\end{enumerate}
Here we say $X \in \Ob(\mathcal{C})$ is connected if
it is not initial and for any monomorphism $Y \to X$
either $Y$ is initial or $Y \to X$ is an isomorphism.
\end{definition}

\begin{example}
\label{example-from-C-F-to-G-sets}
Let $\mathcal{C}$ be a category and let $F : \mathcal{C} \to \textit{Sets}$
be a functor such that $F(X)$ is finite for all $X \in \Ob(\mathcal{C})$.
By Lemma \href{https://stacks.math.columbia.edu/tag/0BMR}{lemma aut inverse limit (Tag 0BMR)} we see that $G = \text{Aut}(F)$
comes endowed with the structure of a profinite topological group in a
canonical manner. We obtain a functor
\begin{equation}
\label{equation-remember}
\mathcal{C} \longrightarrow \textit{Finite-}G\textit{-Sets},\quad
X \longmapsto F(X)
\end{equation}
where $F(X)$ is endowed with the induced action of $G$. This action
is continuous by our construction of the topology on $\text{Aut}(F)$.
\end{example}

\begin{proposition}
\label{proposition-galois}
\begin{reference}
This is a weak version of \href{https://stacks.math.columbia.edu/bibliography}{Bibliography: SGA1 (Expos\textbackslash{}'e V)}.
The proof is borrowed from \href{https://stacks.math.columbia.edu/bibliography}{Bibliography: BS (Theorem 7.2.5)}.
\end{reference}
Let $(\mathcal{C}, F)$ be a Galois category. Let $G = \text{Aut}(F)$
be as in Example \ref{example-from-C-F-to-G-sets}. The functor
$F : \mathcal{C} \to \textit{Finite-}G\textit{-Sets}$
(\ref{equation-remember}) an equivalence.
\end{proposition}

\begin{proof}
We will use the results of Lemma \href{https://stacks.math.columbia.edu/tag/0BN0}{lemma epi mono (Tag 0BN0)} without further mention.
In particular we know the functor is faithful.
By Lemma \href{https://stacks.math.columbia.edu/tag/0BN3}{lemma tame (Tag 0BN3)} we know that for any connected $X$ the
action of $G$ on $F(X)$ is transitive. Hence for general $X$ the functor
$F$ sends the connected components of $X$ (which exist by
an axiom of a Galois category) to the $G$-orbits of $F(X)$.
Let $X$ and $Y$ be objects and let
$s : F(X) \to F(Y)$ be a map. Then the graph
$\Gamma_s \subset F(X) \times F(Y)$ of $s$
is a union of connected components. Hence there exists a
union of connected components $Z$ of $X \times Y$,
which comes equipped with a monomorphism $Z \to X \times Y$,
with $F(Z) = \Gamma_s$. Since $F(Z) \to F(X)$ is bijective
we see that $Z \to X$ is an isomorphism and we conclude
that $s = F(f)$ where $f : X \cong Z \to Y$ is the composition.
Hence $F$ is fully faithful.

\medskip\noindent
To finish the proof we show that $F$ is essentially surjective.
It suffices to show that $G/H$ is in the essential image for
any open subgroup $H \subset G$ of finite index.
By definition of the topology on $G$ there exists a finite
collection of objects $X_i$ such that
$$
\Ker(G \longrightarrow \prod\nolimits_i \text{Aut}(F(X_i)))
$$
is contained in $H$. We may assume $X_i$ is connected
for all $i$. We can choose a Galois object $Y$ mapping
to a connected component of $\prod X_i$ using
Lemma \href{https://stacks.math.columbia.edu/tag/0BN2}{lemma galois (Tag 0BN2)}. Choose an isomorphism $F(Y) = G/U$
in $G\textit{-sets}$ for some open subgroup $U \subset G$.
As $Y$ is Galois, the group
$\text{Aut}(Y) = \text{Aut}_{G\textit{-Sets}}(G/U)$ acts transitively
on $F(Y) = G/U$. This implies that $U$ is normal. Since
$F(Y)$ surjects onto $F(X_i)$ for each $i$ we see that
$U \subset H$. Let $M \subset \text{Aut}(Y)$ be the finite subgroup
corresponding to
$$
(H/U)^{opp} \subset (G/U)^{opp} = \text{Aut}_{G\textit{-Sets}}(G/U)
= \text{Aut}(Y).
$$
Set $X = Y/M$, i.e., $X$ is the coequalizer
of the arrows $m : Y \to Y$, $m \in M$.
Since $F$ is exact we see that $F(X) = G/H$ and the
proof is complete.
\end{proof}
\end{document}
